\setcounter{section}{3}
\section{Phương trình quy về phương trình bậc nhất một ẩn}

\begin{lesson}
    \subsection{Kiến thức trọng tâm}

    \begin{center}
        \begin{tblr}{
            width=\linewidth,
            colspec={X[l] X[l]},
            cells={valign=t},
            row{1}={font=\bfseries, halign=c}
        }
            Khái niệm, thuật ngữ & Kiến thức, kĩ năng \\
            \textbullet\ Phương trình tích\par
            \textbullet\ Phương trình chứa ẩn ở mẫu
            &
            \textbullet\ Giải phương trình tích có dạng $(ax+b)(cx+d)=0$.\par
            \textbullet\ Giải phương trình chứa ẩn ở mẫu quy về phương trình bậc nhất.
        \end{tblr}
    \end{center}

    Trong một khu vườn hình vuông có cạnh bằng 15 m người ta làm một lối đi xung quanh vườn có bề rộng là $x$ (m) (H.2.1). Để diện tích phần đất còn lại là $169\,\mathrm{m}^2$ thì bề rộng $x$ của lối đi là bao nhiêu?

    \begin{center}
        \begin{tikzpicture}
            \coordinate (A) at (0,0);
            \coordinate (B) at (4.8,0);
            \coordinate (C) at (4.8,4.8);
            \coordinate (D) at (0,4.8);
            \coordinate (E) at (0.7,0.7);
            \coordinate (F) at (4.1,0.7);
            \coordinate (G) at (4.1,4.1);
            \coordinate (H) at (0.7,4.1);
            \draw (A)--(B)--(C)--(D)--cycle;
            \draw (E)--(F)--(G)--(H)--cycle;
            \draw[<->] (D) -- node[above] {$15\,\mathrm{m}$} (C);
            \draw[<->] (G) -- node[above] {$x$} (C |- G);
        \end{tikzpicture}

        \textit{Hình 2.1}
    \end{center}

    \textbf{1. PHƯƠNG TRÌNH TÍCH}

    \textbf{Phương trình tích}

    \textbf{HĐ1} Phân tích đa thức $P(x)=(x+1)(2x-1)+(x+1)x$ thành nhân tử.

    \answerline
    \answerline

    \textbf{HĐ2} Giải phương trình $P(x)=0$.

    \answerline
    \answerline

    \begin{keyconcept}
        Để giải phương trình tích $(ax+b)(cx+d)=0$, ta giải hai phương trình $ax+b=0$ và $cx+d=0$. Sau đó lấy tất cả các nghiệm của chúng.
    \end{keyconcept}

    \textbf{Ví dụ 1:} Giải phương trình $(2x+1)(3x-1)=0$.

    \textbf{Giải}

    Ta có $(2x+1)(3x-1)=0$

    nên $2x+1=0$ hoặc $3x-1=0$.
    \begin{itemize}
        \item $2x+1=0$ hay $2x=-1$, suy ra $x=-\dfrac{1}{2}$.
        \item $3x-1=0$ hay $3x=1$, suy ra $x=\dfrac{1}{3}$.
    \end{itemize}
    Vậy phương trình đã cho có hai nghiệm là $x=-\dfrac{1}{2}$ và $x=\dfrac{1}{3}$.

    \textbf{Ví dụ 2:} Giải phương trình $x^2-x=-2x+2$.

    \textbf{Giải}

    Biến đổi phương trình đã cho về phương trình tích như sau:
    \[
        \begin{aligned}
            x^2-x&=-2x+2\\
            x^2-x+2x-2&=0\\
            x(x-1)+2(x-1)&=0\\
            (x+2)(x-1)&=0.
        \end{aligned}
    \]
    Ta giải hai phương trình sau:
    \begin{itemize}
        \item $x+2=0$ suy ra $x=-2$.
        \item $x-1=0$ suy ra $x=1$.
    \end{itemize}
    Vậy phương trình đã cho có hai nghiệm là $x=-2$ và $x=1$.

    \textbf{Nhận xét.} Trong Ví dụ 2, ta thực hiện việc giải phương trình theo hai bước:
    \begin{itemize}
        \item \textbf{Bước 1.} Đưa phương trình về phương trình tích $(ax+b)(cx+d)=0$;
        \item \textbf{Bước 2.} Giải phương trình tích tìm được.
    \end{itemize}

    \textbf{Luyện tập 1.} Giải các phương trình sau:
    \begin{enumerate}[label=\alph*)]
        \item $(3x+1)(2-4x)=0$;

        \answerline
        \answerline
        \item $x^2-3x=2x-6$.

        \answerline
        \answerline
    \end{enumerate}

    \textbf{Vận dụng.} Giải bài toán ở tình huống mở đầu.

    \answerline
    \answerline
    \answerline

    \textbf{2. PHƯƠNG TRÌNH CHỨA ẨN Ở MẪU}

    \textbf{Điều kiện xác định của một phương trình}

    Xét phương trình
    \[
        x+\dfrac{1}{x+1}=-1+\dfrac{1}{x+1}. \tag{1}
    \]

    \textbf{HĐ3} Chuyển các biểu thức chứa ẩn từ vế phải sang vế trái, rồi thu gọn vế trái.

    \answerline
    \answerline

    \textbf{HĐ4} Giá trị $x=-1$ có là nghiệm của phương trình đã cho hay không? Vì sao?

    \answerline
    \answerline

    Khi giải phương trình chứa ẩn ở mẫu, ta phải chú ý đến điều kiện xác định của một phương trình. Chẳng hạn, đối với phương trình trên, để phương trình xác định thì $x+1\ne0$ hay $x\ne-1$.

    \begin{keyconcept}
        Đối với phương trình chứa ẩn ở mẫu, ta thường đặt điều kiện cho ẩn để tất cả các mẫu thức trong phương trình đều khác 0 và gọi đó là điều kiện xác định (viết tắt là ĐKXĐ) của phương trình.
    \end{keyconcept}

    \textbf{Ví dụ 3:} Tìm điều kiện xác định của mỗi phương trình sau:
    \begin{enumerate}[label=\alph*)]
        \item $\dfrac{5x+2}{x-1}=0$;
        \item $\dfrac{1}{x+1}=1+\dfrac{1}{x-2}$.
    \end{enumerate}

    \textbf{Giải}

    a) Vì $x-1\ne0$ khi $x\ne1$ nên ĐKXĐ của phương trình $\dfrac{5x+2}{x-1}=0$ là $x\ne1$.

    b) Vì $x+1\ne0$ khi $x\ne-1$ và $x-2\ne0$ khi $x\ne2$ nên ĐKXĐ của phương trình $\dfrac{1}{x+1}=1+\dfrac{1}{x-2}$ là $x\ne-1$ và $x\ne2$.

    \textbf{Luyện tập 2.} Tìm điều kiện xác định của mỗi phương trình sau:
    \begin{enumerate}[label=\alph*)]
        \item $\dfrac{3x+1}{2x-1}=1$;

        \answerline
        \item $\dfrac{x}{x-1}+\dfrac{x+1}{x}=2$.

        \answerline
    \end{enumerate}

    \textbf{Cách giải phương trình chứa ẩn ở mẫu}

    \textbf{HĐ5} Xét phương trình
    \[
        \dfrac{x+3}{x}=\dfrac{x+9}{x-3}. \tag{2}
    \]
    Hãy thực hiện các yêu cầu sau để giải phương trình (2):
    \begin{enumerate}[label=\alph*)]
        \item Tìm điều kiện xác định của phương trình (2);

        \answerline
        \item Quy đồng mẫu hai vế của phương trình (2), rồi khử mẫu;

        \answerline
        \item Giải phương trình vừa tìm được;

        \answerline
        \item Kết luận nghiệm của phương trình (2).

        \answerline
    \end{enumerate}

    \begin{keyconcept}
        Để giải phương trình chứa ẩn ở mẫu ta thường thực hiện các bước như sau:
        \begin{itemize}
            \item \textbf{Bước 1.} Tìm điều kiện xác định của phương trình.
            \item \textbf{Bước 2.} Quy đồng mẫu hai vế của phương trình rồi khử mẫu.
            \item \textbf{Bước 3.} Giải phương trình vừa tìm được.
            \item \textbf{Bước 4 (Kết luận).} Trong các giá trị tìm được của ẩn ở Bước 3, giá trị nào thoả mãn điều kiện xác định chính là nghiệm của phương trình đã cho.
        \end{itemize}
    \end{keyconcept}

    \textbf{Ví dụ 4:} Giải phương trình
    \[
        \dfrac{2}{x+1}+\dfrac{1}{x-2}=\dfrac{3}{(x+1)(x-2)}. \tag{3}
    \]

    \textbf{Giải}

    Điều kiện xác định $x\ne-1$ và $x\ne2$.

    Quy đồng mẫu và khử mẫu, ta được
    \[
        \dfrac{2(x-2)}{(x+1)(x-2)}+\dfrac{x+1}{(x+1)(x-2)}
        =\dfrac{3}{(x+1)(x-2)},
    \]
    suy ra
    \[
        2(x-2)+(x+1)=3. \tag{3a}
    \]
    Giải phương trình (3a):
    \[
        \begin{aligned}
            2(x-2)+(x+1)&=3\\
            2x-4+x+1&=3\\
            3x-3&=3\\
            3x&=6\\
            x&=2.
        \end{aligned}
    \]
    Giá trị $x=2$ không thỏa mãn ĐKXĐ. Vậy phương trình (3) vô nghiệm.

    \textbf{Luyện tập 3.} Giải phương trình
    \[
        \dfrac{1}{x-1}-\dfrac{4x}{x^3-1}=\dfrac{x}{x^2+x+1}.
    \]

    \answerline
    \answerline
    \answerline
\end{lesson}

\begin{exercises}
    \subsection{Bài tập}

    \begin{questions}[series=SGK]
        \item Giải các phương trình sau:
        \begin{questions}
            \item $x(x-2)=0$;

            \answerline
            \answerline
            \item $(2x+1)(3x-2)=0$.

            \answerline
            \answerline
        \end{questions}

        \item Giải các phương trình sau:
        \begin{questions}
            \item $(x^2-4)+x(x-2)=0$;

            \answerline
            \answerline
            \item $(2x+1)^2-9x^2=0$.

            \answerline
            \answerline
        \end{questions}

        \item Giải các phương trình sau:
        \begin{questions}
            \item $\dfrac{2}{2x+1}+\dfrac{1}{x+1}=\dfrac{3}{(2x+1)(x+1)}$;

            \answerline
            \answerline
            \answerline
            \item $\dfrac{1}{x+1}-\dfrac{x}{x^2-x+1}=\dfrac{3x}{x^3+1}$.

            \answerline
            \answerline
            \answerline
        \end{questions}

        \item Bác An có một mảnh đất hình chữ nhật với chiều dài 14 m và chiều rộng 12 m. Bác dự định xây nhà trên mảnh đất đó và dành một phần diện tích đất để làm sân vườn như Hình 2.3. Biết diện tích đất làm nhà là $100\,\mathrm{m}^2$. Hỏi $x$ bằng bao nhiêu mét?

        \begin{center}
            \begin{tikzpicture}
                \coordinate (A) at (0,0);
                \coordinate (B) at (5.6,0);
                \coordinate (C) at (5.6,4.8);
                \coordinate (D) at (0,4.8);
                \coordinate (E) at (4.0,0);
                \coordinate (F) at (4.0,4.0);
                \coordinate (G) at (0,4.0);
                \draw (A)--(B)--(C)--(D)--cycle;
                \draw (A)--(E)--(F)--(G)--cycle;
                \node at (2.0,2.0) {Nhà};
                \node at (3.8,4.4) {Sân vườn};
                \draw[<->] (D) -- node[above] {$14\,\mathrm{m}$} (C);
                \draw[<->] (A) -- node[left] {$12\,\mathrm{m}$} (D);
                \draw[<->] (G) -- node[left] {$x$} (D);
                \draw[<->] (E) -- node[below] {$x+2$} (B);
            \end{tikzpicture}

            \textit{Hình 2.3}
        \end{center}

        \answerline
        \answerline
        \answerline
        \answerline

        \item Hai người cùng làm chung một công việc thì xong trong 8 giờ. Hai người cùng làm được 4 giờ thì người thứ nhất bị điều đi làm công việc khác. Người thứ hai tiếp tục làm việc trong 12 giờ nữa thì xong công việc. Gọi $x$ là thời gian người thứ nhất làm một mình xong công việc (đơn vị tính là giờ, $x>0$).
        \begin{questions}
            \item Hãy biểu thị theo $x$:
            \begin{itemize}
                \item Khối lượng công việc mà người thứ nhất làm được trong 1 giờ;
                \item Khối lượng công việc mà người thứ hai làm được trong 1 giờ.
            \end{itemize}

            \answerline
            \answerline
            \item Hãy lập phương trình theo $x$ và giải phương trình đó. Sau đó cho biết, nếu làm một mình thì mỗi người phải làm trong bao lâu mới xong công việc đó.

            \answerline
            \answerline
            \answerline
            \answerline
        \end{questions}
    \end{questions}
\end{exercises}

\begin{practice}
    \subsection{Luyện tập}

    \begin{questions}[series=SBT]
        \item Giải các phương trình sau:
        \begin{questions}
            \item $(x+2)^2-(2x+1)(x+2)=0$;

            \answerline
            \answerline
            \item $16x^2-(3x+2)^2=0$.

            \answerline
            \answerline
        \end{questions}

        \item Giải các phương trình sau:
        \begin{questions}
            \item $x^3+3x^2-8=x^3+2x^2-7$;

            \answerline
            \answerline
            \item $x(2x-5)=(2x+1)(5-2x)$.

            \answerline
            \answerline
        \end{questions}

        \item Giải các phương trình sau:
        \begin{questions}
            \item $x^2+x=-6x-6$;

            \answerline
            \answerline
            \item $2x^2-2x=x-1$.

            \answerline
            \answerline
        \end{questions}

        \item Giải các phương trình sau:
        \begin{questions}
            \item $\dfrac{5x-1}{3x+2}-\dfrac{5x+2}{3x}=0$;

            \answerline
            \answerline
            \answerline
            \item $\dfrac{6x-5}{2x-1}-\dfrac{9x}{3x-1}=0$.

            \answerline
            \answerline
            \answerline
        \end{questions}

        \item Giải các phương trình sau:
        \begin{questions}
            \item $\dfrac{3}{x+2}+\dfrac{x}{x^2-2x+4}=\dfrac{4x^2}{x^3+8}$;

            \answerline
            \answerline
            \answerline
            \item $\dfrac{3}{2x+1}+\dfrac{7}{3x+2}=\dfrac{21x+10}{(2x+1)(3x+2)}$.

            \answerline
            \answerline
            \answerline
        \end{questions}

        \item Một vật rơi tự do từ độ cao so với mặt đất là 120 mét. Bỏ qua sức cản không khí, quãng đường chuyển động $s$ (mét) của vật rơi tự do sau thời gian $t$ được biểu diễn gần đúng bởi công thức $s=4{,}9t^2$, trong đó $t$ là thời gian tính bằng giây. Sau bao nhiêu giây kể từ khi bắt đầu rơi thì vật này chạm mặt đất (làm tròn kết quả đến chữ số hàng đơn vị)?

        \answerline
        \answerline
        \answerline
    \end{questions}
\end{practice}
